\title[SNode.C]{\large A \alert{simple}, lightweight, highly extensible, event driven, layer-based framework for network applications in the spirit of \alert{node}.js written entirely in \alert{C}++}

\def\insertlecturepartnumbername{\Huge\alert{SNode.C}}

\mode<presentation>{\lecture[Überblick]{Überblick}{Überblick}}

\part{Einordnung}
\section{Motivation}

\begin{frame}{Motivation}{Einordnung}
  \begin{center}
    \alert{\Huge \color{fh-ooe-red} SNode.C}\\\scriptsize\copyright{} 2020, 2021, 2022, 2023 Volker Christian\\
    \href{https://github.com/VolkerChristian/snode.c}{\beamergotobutton{https://github.com/VolkerChristian/snode.c}} \href{https://www.gnu.org/licenses/lgpl-3.0.en.html}{\beamergotobutton{GNU Lesser General Public License v3.0}}
  \end{center}
  Die Entwicklung von \texttt{SNode.C} begann w\"ahrend des ersten Corona-Lockdowns im Sommersemester 2020 im Rahmen der Master-Lehrveranstaltung \alert{Networked and
    Distributed Systems}.
  \begin{description}[\alert{Designentscheidung}]
    \item [\alert{Ziel}] Studierende mit den grundlegenden \alert{Techniken und Konzepten von Netzwerk- und Webframeworks} vertraut zu machen.
    \item [\alert{Vorbild}] Bekanntes \alert{express} Framework von \alert{node.js}.
    \item [\alert{Designentscheidung}] ursprünglich
    \begin{itemize}
      \item Single tasking, single threaded.
      \item Non blocking, event driven, multiplexed I/O.
      \item "`select"' als einziger I/O-Multiplexer
      \item API so \"ahnlich wie m\"oglich (JS vs. C++) dem API von node.js-express.
      \item Nur \alert{serverseitig}.
      \item Nur \alert{IPv4}
    \end{itemize}
    \item [\alert{Aktueller Stand}] Viel mehr als anfangs geplant
  \end{description}
\end{frame}
